%%%%

%%%   Самарский государственный технический университет

%%%   Кафедра прикладной математики и информатики

%%%

%%%   Преамбула для статьи в журнал "Вестник Самарского государственного

%%%   технического университета. Серия: «Физико-математические науки»"

%%%   Ориентирована под MikTEx 2.4 for Win32

%%%

%%%   Хотя сам журнал собирается под GNU/Linux на дистрибутиве TeXLive



\documentclass[11pt,a4paper,twoside]{article}



\usepackage[cp1251]{inputenc}

\usepackage[T2A]{fontenc}

\usepackage[english,russian]{babel}

\usepackage{samgtubib}

\usepackage[dvips]{graphicx}

\usepackage[multiple]{footmisc}



\usepackage{samgtu}

\renewcommand{\VestnikYear}{2009} % Год издания Вестника

\renewcommand{\VestnikNumber}{2\,(19)} % Номер Вестника

\renewcommand{\VestnikMonth}{Ноябрь} % Месяц выхода Вестника

\renewcommand{\VestnikData}{01.11.09} % Дата подписания Вестника в печать

\renewcommand{\VestnikUPL}{26,64} % Усл. печ. л.

\renewcommand{\VestnikUIL}{25,73} % Уч.-изд. л.

\renewcommand{\VestnikRegNumber}{479/09} % Регистрационный номер

\renewcommand{\VestnikOrder}{994} % Номер заказа







%

%\begin{document}







\renewcommand{\firstpage}{116}

\renewcommand{\lastpage}{122}







\udk{539.37}

\msc{74K10; 74E10, 74B05, 35J15}

%*74K10 Rods (beams, columns, shafts, arches, rings, etc.)

%74E10 Anisotropy

%*74E10 Anisotropy

%74B05 Classical linear elasticity

%*35J15 Second order elliptic equations, general



\title{Антиплоская деформация цилиндрически анизотропного упругого стержня}{Антиплоская деформация цилиндрически анизотропного стержня \dots}



\author{Ю.\,А.~Боган }{Б\,о\,г\,а\,н\,~Ю.\,А.}



\institute{\igil.}



\email{E-mail}{bogan@hydro.nsc.ru}



\otherinf{\fullauthor{Юрий Александрович Боган}\regalia{д.ф.-м.н.} \position{старший научный сотрудник} \dept{отд. механики деформируемого твёрдого тела}}



\keywords{цилиндрическая анизотропия, упругость}



\workabstract{Построены явные решения краевых задач Дирихле и~Неймана при антиплоской деформации цилиндрически анизотропного

материала для круговой области. Доказана однозначная разрешимость задачи Дирихле в ограниченной области на плоскости с~кусочно"=гладкой границей. }





\engtitle{Antiplane deformation of a cylindrically anisotropic elastic rod}







\engauthor{Yu.\,A.~Bogan}{B\,o\,g\,a\,n~Yu.\,A. }





\enginstitute{\engigil.}



\otherenginf{\fullauthor{Yurii A. Bogan} \regalia{Dr.\,Sci. (Phys. {\rm \&} Math.)} 

\position{Senior Research Scientist} \dept{Dept. of Deformable Solid Body}}



\engabstract{The problem of antiplane deformation of general cylindrical anisotropic material is studied in this paper. Explicit solutions of Dirichlet and Neumann problems are given for a circular domain. The existence of unique weak solution of the Dirichlet problem in a bounded region with a piece-wise smooth boundary is proved.}



\engkeywords{cylindrical anisotropy, elasticity}



\date{27/I/2012}

\revisiondate{27/II/2012}



%%%%%%%%%%%%%%%%%%%%%%%%%%%%%%%%%%%%%%%%%%%%%%%%%%%%%%%%%%%%%%%%%%%%%%%%%%%%%%%%%%%%





\maketitle





\Section[n]{Введение} Краевые задачи кручения упругих анизотропных

стержней хорошо изучены в монографии С.\,Г.~Лехницкого

\cite{bogan:bib1}. Некоторые результаты в этом направлении были

получены  Тингом в статье \cite{bogan:bib2}. Одна задача о щели

рассматривалась в статье \cite{bogan:bib3}. Необходимо отметить,

что в отмеченных работах рассматривался только цилиндрически

ортотропный материал. Здесь же изучены краевые задачи Дирихле и

Неймана для произвольного цилиндрически анизотропного материала в

предположении антиплоской деформации. При произвольной анизотропии

в это уравнение входит смешанная производная по $ r$, $\theta $,

что существенно усложняет построение явного решения. Аналогичное

уравнение возникает в стационарной теплопроводности

\cite{bogan:bib4}.



\smallskip



\Section{Основные уравнения и постановка задачи} Введем

цилиндрическую систему координат $ r$, $\theta$, $z$. Тогда в

предположении антиплоской деформации перемещения даются формулами

$$

\begin{array}{l}

\displaystyle u_r= z\Bigr(-\dfrac{1}{r}\dfrac{\partial w}{\partial \theta} +

a_{45} \tau_{\theta z} + a_{55} \tau_{rz} \Bigr) + c_1 \cos

\theta +c_2 \sin \theta,\\

\displaystyle u_\theta= z \Bigl( - \dfrac{\partial w}{\partial r} + a_{44}

\tau_{\theta z} + a_{45} \tau_{rz} \Bigr) + c_2 \cos \theta -c_1

\sin \theta + c_0 r.

\end{array}

$$

 Здесь $c_0$, $c_1$, $c_2$ "--- вещественные

постоянные. Вторые слагаемые в предыдущих формулах отвечают за

жёсткое перемещение тела. Фиксируем их каким"=либо способом. 

В~предположении независимости перемещений  от $z$  напряжения определяются из соотношений

\begin{gather}

\label{bogan:eq1}

- \dfrac{1}{r}\dfrac{\partial w}{\partial

\theta} + a_{45} \tau_{\theta z} + a_{55} \tau_{rz} =0,\quad 

 - \dfrac{\partial w}{\partial r} + a_{44} \tau_{\theta z} +

a_{45} \tau_{rz} =0.

\end{gather}



Решив  уравнения \eqref{bogan:eq1} относительно напряжений, получим

\begin{gather*}

\tau_{\theta z}= A_{44}\frac{\partial w}{r \partial

\theta}+A_{45}\frac{\partial w}{\partial r},\quad 

\tau_{r z}=A_{45}\frac{\partial w}{r \partial

\theta}+A_{55}\frac{\partial w}{\partial r},

\end{gather*}

где

$A_{44}={a_{55}}/{\delta}$, $A_{45}=A_{54}=-{a_{45}}/{\delta}$, $A_{55}={a_{44}}/{\delta}$, $\delta=a_{44}a_{55}-a_{45}^2

>0$.

Здесь $ A_{44}$, $ A_{45}$, $ A_{55}$ "--- функции координат. 

При подстановке результата в единственное уравнение равновесия

\begin{equation*}

\frac{\partial \tau_{r z} }{\partial r}+ \frac{1}{r}\frac{\partial

\tau_{\theta z}}{\partial \theta }+\frac{\tau_{r

z}}{r}=f(r,\theta)

\end{equation*}

получим, что $w(r,\,\theta)$ определяется из уравнения

\begin{multline}

\label{bogan:eq2}

\frac{\partial}{\partial r}\Bigl( A_{55}\frac{\partial w}{\partial r}+

A_{45}\frac{1}{r}\frac{\partial w}{\partial \theta}\Bigr )+

\frac{1}{r}\frac{\partial}{\partial \theta}\Bigl(A_{45}\frac{\partial w}{\partial r}+

A_{44}\frac{1}{r}\frac{\partial w}{\partial \theta}\Bigr) +\\

+\frac{1}{r}\Bigl( A_{55}\frac{\partial w}{\partial

r}+A_{45}\frac{1}{r}\frac{\partial w}{\partial \theta}\Bigr

)=f(r, \theta).

\end{multline}

В частности, для однородного стержня при постоянных модулях сдвига и равной нулю правой части имеем

\begin{equation}

\label{bogan:eq3}

 A_{55}\biggl(\frac{\partial^2 w}{\partial

r^2}+\frac{1}{r}\frac{\partial w}{\partial

r}\biggr)+2A_{45}\frac{1}{r}\frac{\partial^2 w}{{\partial

r}{\partial \theta}}+ A_{44}\frac{1}{r^2}\frac{\partial^2

w}{\partial \theta^2}=0.

\end{equation}

С уравнением \eqref{bogan:eq2} можно ассоциировать билинейную симметричную форму

\begin{multline}

\label{bogan:eq4}

 a(w, \psi)=\int_{\partial Q} 

 \Bigl [ \Bigl(A_{45}\dfrac{1}{r}\dfrac{\partial w}{\partial \theta}+A_{55}\dfrac{\partial w}{\partial r} \Bigr)\dfrac{\partial

\psi}{\partial r}+ \\

+ \Bigl(A_{44}\dfrac{1}{r}\dfrac{\partial w}{\partial \theta}+A_{45}\dfrac{\partial w}{\partial r}

\Bigr)\dfrac{1}{r}\dfrac{\partial \psi}{\partial \theta} \Bigr

] rdrd\theta

\end{multline}

и соответствующую ей квадратичную форму $a( w,\, w)$. 

Уравнение

\eqref{bogan:eq2} можно рассматривать как уравнение Эйлера для функционала $a(w,\, w)$. 

Перейдём в~\eqref{bogan:eq4} к декартовым координатам $ (x_1,\, x_2)$, где $ x_1 =r \cos \theta$, $x_2= r \sin \theta$:

\begin{equation}

\label{bogan:eq5} a(w,\, w)=\int_{\partial Q}

\Bigl[g_{11}\Bigl(\dfrac{\partial w}{\partial x_1}\Bigr)^2+2 g_{12}\dfrac{\partial w}{\partial x_1}\dfrac{\partial w}{\partial

x_2}+g_{22}\Bigl(\dfrac{\partial w}{\partial x_2}\Bigr)^2 \Bigr] dx_1dx_2,

\end{equation}

где

$g_{11}= A_{44}\sin^2 \theta -2A_{45}\cos \theta \sin \theta+A_{55}\cos^2 \theta$, 

$g_{12}= (A_{55}-A_{44})\sin \theta \cos \theta+2 A_{45}(\cos^2 \theta -\sin^2 \theta)$, 

$g_{21}= g_{12}$,

$g_{22}= A_{44}\cos^2\theta+2 A_{45}\cos \theta \sin \theta + A_{55}\sin^2 \theta$.

В декартовых координатах  уравнение (\ref{bogan:eq2})

можно записать так:

\begin{equation}

\label{bogan:eq6}

 \dfrac{\partial }{\partial x_1}\Bigl(g_{11}(x) \dfrac{\partial w}{\partial x_1}{+} g_{12}(x)\dfrac{\partial

w}{\partial x_2} \Bigr)+ \dfrac{\partial }{\partial

x_2}\Bigl(g_{12}(x) \dfrac{\partial w}{\partial x_1}{+}

g_{22}(x)\dfrac{\partial w}{\partial x_2} \Bigr)=f(x_1,x_2).

\end{equation}

Как видно из последних соотношений, коэффициенты уравнения "--- ограниченные измеримые функции, при этом разрывные в начале

координат. 

Действительно,  $ \sin \theta \cos \theta= x_1 x_2/(x_1^2 +x_2^2)$ "--- известный пример разрывной функции.

Предположим, что  правая часть $ f(x_1,\,x_2)\in L^2(Q) $, 

 $Q$ "---   ограниченная область на плоскости $\mathbb R^2$, и  существуют постоянные $C_1$, $C_2 > 0$, такие, что

\[

 C_1 |\xi|^2 \leqslant g_{ij}(x)\xi_i \xi_j \leqslant C_2 |\xi|^2

  \]

для любых $x \in Q$ и всех $ \xi=(\xi_1,\,\xi_2)\in \mathbb R^2$.

Обобщённое решение краевой задачи определяется стандартным образом как функция $ w(x_1,\,x_2) \in H_0^1(Q)$, удовлетворяющая

интегральному тождеству

\begin{equation}

\label{bogan:eq7} \int_Q g_{ij} w_{x_i} \varphi_{x_j} dx=\int_Q f(x)\varphi(x) dx

\end{equation}

для любой $ \varphi(x) \in C_0^\infty(Q) $. 

В~\cite{bogan:bib5} доказано существование единственного слабого решения задачи Дирихле для эллиптического дифференциального оператора в~дивергентной форме в~предположении равномерной эллиптичности и~измеримости коэффициентов. 

Очевидно, что предыдущее уравнение таковым является, так как матрица $ g=(g_{ij})$ получается поворотом из  матрицы $ A=(A_{ij})$ и поэтому имеет те же собственные значения, сохраняя при этом свойство положительной определённости.



Стоит отметить, что подобные уравнения возникают при изучении вопроса о непрерывности по Гёльдеру  решений эллиптических

уравнений второго порядка \cite{bogan:bib6} и служат примерами эллиптических уравнений с~негладкими (в начале координат)

решениями. 

При использовании метода характеристик

\cite{bogan:bib2} и введении новых независимых неизвестных

\[

 \xi= \varrho\cos \Bigl( \dfrac{\alpha}{\beta} \ln \varrho - \theta \Bigr), \quad

 \eta= \varrho \sin \Bigl( \dfrac{\alpha}{\beta} \ln \varrho - \theta \Bigr), \quad \varrho=r^\beta

 \]

исходное  уравнение обращается в уравнение Пуассона

\[

 \dfrac{\partial^2 u(\xi,\,\eta)}{\partial \xi^2}+\dfrac{\partial^2 u(\xi,\, \eta)}{\partial \eta^2}= f(\xi,\, \eta).

  \]

Здесь $u(\xi,\,\eta)=w(r,\, \theta)$. 

При этом семейство характеристик само имеет особую точку; действительно, семейство

кривых $\theta- ({\alpha}/{\beta}) \ln \varrho ={\rm const}$ "--- семейство логарифмических спиралей с особой точкой (фокусом в начале координат). При $\alpha =0$ фокус обращается в

центр, а семейство логарифмических спиралей обращается в семейство

концентрических окружностей. Последнее обстоятельство, в

частности, показывает, что начало координат "--- особая точка и

для уравнения Лапласа. 

Наличие особой точки и приводит, вообще говоря, к сингулярности решения в начале координат. 

Ниже даны явные решения задач Неймана и Дирихле для круговой области с~центром в начале координат.



\smallskip



\Section{Задача Неймана для круга} Построим  в круге радиуса $ R> 0 $ с центром в начале координат формальное решение уравнения

\eqref{bogan:eq3} при граничном условии

\begin{equation}

\label{bogan:eq8}

 \tau_{r z}(R,\,\theta)=A_{45}\frac{\partial w}{r

\partial \theta}+A_{55}\frac{\partial w}{\partial r}=g(\theta),

\end{equation}

где $ g(\theta) $ "--- непрерывная суммируемая  периодическая функция окружной координаты. 

Потребуем ограниченности решения в начале координат. 

Применим  метод разделения переменных. 

Представим $ g(\theta) $ в виде ряда Фурье 

\[

 g(\theta)= \frac{a_0}{2}+ \sum_{n=1}^\infty( a_n\cos n\theta+b_n\sin n\theta),

  \]

где

\[

 a_n=\dfrac{1}{\pi}\int_0^{2\pi} g(\varphi) \cos n \varphi d\varphi,\quad

  b_n= \dfrac{1}{\pi}\int_0^{2\pi} g(\varphi) \cos n \varphi d\varphi, \quad  n=0,\,1, \ldots.

  \]



Ищем решение уравнения \eqref{bogan:eq3} в виде

\[

 w(r,\,\theta)=f_0(r)+ \sum_{n=1}^\infty( f_n(r)\cos n\theta+g_n(r)\sin n\theta),

\]

где 

$ f_n(r)$, $g_n(r) $ определяются из системы

$$

\begin{array}{r}

\displaystyle

 A_{55}\Bigl(\dfrac{d^2 f_n}{dr^2}+\dfrac{1}{r}\dfrac{df_n}{dr} \Bigr)-

 \dfrac{n^2A_{44}}{r^2}f_n+\dfrac{2n A_{45}}{r}\dfrac{dg_n}{dr}=0, \\

 \displaystyle

 -\dfrac{2n A_{45}}{r}\dfrac{df_n}{dr}+ A_{55} \Bigl(\dfrac{d^2 g_n}{dr^2}+\dfrac{1}{r}\dfrac{dg_n}{dr} \Bigr)-\dfrac{n^2A_{44}}{r^2}g_n=0.

 \end{array}

 $$

Положим $ f_n(r)=Br^{\lambda n}$, $g_n(r)=Cr^{\lambda n} $. Тогда

$ B$ и $C$ определяются из 

$$

\phantom{-\!}\left(A_{55}\lambda^2-A_{44}

\right)B+2A_{45}\lambda C=0,\quad

 -2A_{45}\lambda B+ \left(A_{55}\lambda^2-A_{44} \right)C=0,

$$

а $\lambda $ "--- из уравнения

\[

\left(A_{55}\lambda^2-2A_{44} \right)^2+4A_{45}^2\lambda^2=0,

\]

которое допускает представление 

\[

\left(A_{55}\lambda^2-A_{44}-2A_{45}i\lambda

\right)(A_{55}\lambda^2-A_{44}+2A_{45}i\lambda)=0

\]

и  позволяет  найти комплексно сопряженные корни

$\lambda_1=i\alpha+\beta$, $\lambda_3=-i\alpha+\beta$,

$\lambda_2=-i\alpha-\beta$, $\lambda_4=i\alpha-\beta$. Здесь

$\alpha=-{A_{45}}/{A_{55}}$, $\beta={ \sqrt{A_{44}A_{55}-A_{45}^2}}/{A_{55}}$.

Очевидно, $ \beta>0 $, так как предполагается, что $

A_{44}A_{55}-A_{45}^2 >0$. Корни $ \lambda_2$, $\lambda_4 $

следует отбросить, как приводящие к неограниченному росту решения

в начале координат. При $ n=0 $ имеем $ f_0(r)={\rm const}$, а для

функций $ f_n(r)$, $g_n(r) $ получим 

$$

f_n(r)=B_{n1}r^{(i\alpha+\beta)n} +B_{n3}r^{(-i\alpha+\beta)n},\quad 

g_n(r)=-iB_{n1}r^{(i\alpha+\beta)n} +iB_{n3}r^{(-i\alpha+\beta)n}.

$$

Тогда

\[

\tau_{rz}= \sum_{n=0}^\infty \Bigl[ \Bigl( A_{55}\dfrac{d f_n}{dr}+A_{45}n \dfrac{g_n}{r}\Bigr)

 \cos n\theta +  \Bigl( A_{55}\dfrac{d g_n}{dr}-A_{45}n \dfrac{f_n}{r}\Bigr) \sin n\theta \Bigr].

\]

При этом

$$

\begin{array}{l}

A_{55}\dfrac{d f_n}{dr}+A_{45}n \dfrac{g_n}{r}=\beta A_{55}n \left(B_{n1}r^{n \lambda_1  -1}+ B_{n3} r^{n\lambda_3 -1} \right), \\[2mm]

A_{55}\dfrac{d g_n}{dr}-A_{45}n \dfrac{f_n}{r}= \beta A_{55} n \left(-iB_{n1}r^{ n\lambda_1 -1} + B_{n3} r^{ n\lambda_3 -1} \right),

\end{array}

$$

и, определив $ B_{n1}$, $B_{n3} $ из 

$$

\begin{array}{l}

\beta A_{55}n \left( B_{n1}R^{(i\alpha+\beta)n -1} B_{n3}R^{(-i\alpha+\beta)n-1} \right)=a_n,\\

\beta A_{55} n \left( -iB_{n1}R^{(i\alpha+\beta)n-1} +i B_{n3}R^{(-i\alpha+\beta)n-1} \right)=b_n,

\end{array}

$$

получим, что

\begin{gather*}

f_n(r)=\dfrac{R}{2 \beta A_{55} n}\left(\dfrac{r}{R} \right)^{\beta n} \left[a_n\cos

\left(n\alpha \ln \dfrac{r}{R} \right)- b_n\sin \left(n\alpha \ln \dfrac{r}{R} \right)\right],\\

g_n(r)= \dfrac{R}{2\beta A_{55} n} \left(\dfrac{r}{R}\right)^{\beta n} \left[ a_n\sin \left(n\alpha \ln \dfrac{r}{R} \right)

+b_n\cos \left(n\alpha \ln \dfrac{r}{R}\right) \right].

\end{gather*}

Положим

$w(r,\theta)=u(\rho,\theta)$, $\rho=({r}/{R} )^\beta.

$

В результате решение представляется в виде ряда

\[

w(r,\,\theta)=f_0+ \sum_{n=1}^\infty \dfrac{R}{\beta A_{55} n}\left[a_n\cos n \Bigl(\theta-

\dfrac{\alpha}{\beta}\ln\rho \Bigr) +b_n\sin n \Bigl(\theta-\dfrac{\alpha}{\beta}\ln\rho \Bigr) \right],

\]

который можно просуммировать. Действительно, положим $\delta =\beta A_{55}$

и внесём в предыдущую формулу представление  $a_n$, $b_n$, $n=0$, $1$, $\ldots $, в виде коэффициентов Фурье:

\[

w(r,\,\theta) = \dfrac{R}{\pi \delta} \int_{0}^{2\pi} \sum_{n=0}^\infty g(\varphi) \rho^n

\cos n\Bigl( \varphi -\theta + \dfrac{\alpha}{\beta}\ln \rho\Bigr) d\varphi.

\]



Положим $z=r^\beta e^{i \varphi}$, $t= R^\beta e^{i \psi}$, $\psi= \theta -\alpha \ln(r/R)$.

 Имеет место равенство

\[

 \sum_{n=1}^\infty \dfrac{1}{n}\left(\dfrac{z}{t} \right)^n = - \ln \left(1- \dfrac{z}{t} \right).

\]

Оно позволяет представить $w(r,\,\theta)$ как действительную часть выражения

\begin{equation}

\label{bogan:eq11}

w(r,\,\theta)= \dfrac{R}{\pi \delta}\int_{0}^{2 \pi}g(\varphi)\ln \left(1-\dfrac{z}{t}\right)d\varphi,

\end{equation}

то есть записать в виде, аналогичному решению задачи Неймана для уравнения Лапласа:

\begin{equation}

\label{bogan:eq12}

w(r,\,\theta) = - \dfrac{R}{\pi \delta} \int_{0}^{2\pi} g(\varphi)

\ln \left(1+\rho^2 -2\rho \cos\Bigl(\varphi -\theta + \dfrac{\alpha}{\beta} \ln \rho\Bigr) \right)d\varphi + a_0,

\end{equation}

где $a_0$ "--- произвольная постоянная. Очевидно, что $g(\varphi)$ должна удовлетворять необходимому (и достаточному) условию разрешимости

\[

 \int_0^{2\pi} g(\varphi)d\varphi=0.

\]



\smallskip



\Section{Задача Дирихле} Совершенно аналогично можно построить решение задачи

Дирихле для уравнения (\ref{bogan:eq3}) при равной нулю правой части и граничном условии

$w(R,\,\theta)= f(\theta).

$



Используя выражения для коэффициентов Фурье, $ w(\rho,\,\theta)$ можно записать в виде интеграла:

\[

 w(\rho, \,\theta)=\dfrac{1}{2 \pi}\int _0^{2 \pi} g(\varphi)\biggl(1+2 \sum_{n=1}^\infty

  \rho^n\cos n(\varphi-\alpha \ln \rho +\theta)\biggr) d\theta.

\]

Положим $\varphi +\theta -({\alpha}/{\beta}) \ln \rho= \omega$ и учтём, что

\[

1+2 \sum_{n=1}^\infty \rho^n \cos n \omega= \dfrac{1-\rho^2}{1+\rho^2 -2\rho \cos \omega}.

\]

Тогда $ w(\rho,\,\theta)$ можно записать  в виде аналога интеграла Пуассона для гармонической функции:

\begin{equation}

\label{bogan:eq14}

w(\rho,\,\theta)= \frac{1}{2 \pi}\int \limits_0^{2 \pi} f(\psi)\frac{1-\rho^2}{1-2\rho \cos(\psi-\varphi)+\rho^2}d\psi.

\end{equation}

Здесь $\varphi=\theta-\alpha \ln \rho $. В действительности это настоящий интеграл Пуассона для функции

$u(\xi,\,\eta)= w(r,\,\theta)$,  имеющий смысл для любой непрерывной функции, так как функция

$u(\xi,\,\eta)$ "--- гармоническая в этих переменных. Поскольку

$ \vert \rho \cos (\psi-\varphi)\vert \leq \rho$, можно перейти к пределу под знаком интеграла и получить теорему о среднем:

\[

u(0,\,\theta)= w(0,\,\theta)=\dfrac{1}{2 \pi} \int_0^{2 \pi}f(\theta)d \theta.

\]



Рассмотрим  пример. Пусть $ w(R,\,\theta)= k \cos \theta $. Легко вычислить интеграл Пуассона и увидеть,

что $ w(r,\,\theta)= k \rho \cos (\theta -\gamma \ln \rho) $. Однако при этом функция $ w(r,\,\theta) $

имеет ограниченную гладкость по $r$ при любом $\beta > 0$.

Гладкость решения существенно зависит от величины $\beta $. Действительно, $ w(r,\,\theta)$

непрерывна в начале координат, но ее первая производная по $ r $ неограниченна, если $ \beta < 1 $

и потому напряжения в начале координат неограниченны. Если  $\beta \geqslant 2 $,

то вторые производные решения растут в начале координат, но напряжения уже ограниченны.



В качестве еще одного примера  рассмотрим задачу Дирихле для сектора кругового кольца:

\begin{gather*}

 w(r,\,\alpha \pi)=0, \quad w(r,\,0)=0, \quad w(R,\,\theta)=g(\theta),\\

0<\alpha<2, \quad  g(0)=g(\alpha \pi)=0

\end{gather*}

в предположении, что $ \alpha =0 $.

Её решение можно представить в виде 

\[

w(r,\,\theta)=\sum_{n=1}^{\infty} a_n

\rho ^{n/\alpha}

\sin {n\theta}/{\alpha},

\]

где $a_n$ "--- коэффициенты в разложении в ряд Фурье функции $g(\theta)$

по ортогональной системе функций $\sin{n\theta}/{\alpha}$:

\[

a_n=\frac{2}{\alpha \pi}\int_0^{\alpha \pi} g(\theta)\sin \frac{n

\theta}{\alpha}d\theta.

\]

Гладкость решения зависит от отношения $\beta$, $\theta$ и функции $g(\theta)$.

Очевидно, решение непрерывно в замкнутой области. Оно даже аналитично, если \mbox{$\beta/ \alpha\in \mathbb Z$.} Имеем

\[

\frac{\partial w}{\partial r}(r,\,\theta)=\sum\limits_{n=1}^\infty

\frac{n \beta}{\alpha}\frac{r^{n \beta/\alpha

-1}}{R^{n\beta / \alpha}}a_n \sin\frac{n\theta}{\alpha}.

\]

Ясно, что гладкость радиальной производной зависит от свойств отношения $\beta / \alpha$. Например,

если $\alpha=2$, диск имеет разрез от начала координат до окружности, слагаемое с $n=1$ имеет порядок

$r^{\beta/2-1}$ и если $ \beta \geqslant 2 $, особенности нет.





\begin{thebibliography}{8}



\RBibitem{bogan:bib1}

\by Лехницкий~С.\,Г.

\book Кручение анизотропных и неоднородных стержней

\publaddr М.

\publ Наука

\yr 1971

\totalpages 240

\misctransl

\by Lekhnitskii~S.\,G.

\book Torsion of Anisotropic and Inhomogeneous Rods

\publaddr Moscow

\publ Nauka

\yr 1971

\totalpages 240



\Bibitem{bogan:bib2}

\by Ting~T.\,C.

\paper The remarkable nature of cylindrically anisotropic elastic materials exemplified by an anti-plane deformation

\jour J. Elasticity

\vol 49

\issue 3

\yr 1998

\pages 269--284



\Bibitem{bogan:bib3}

\by Li~B., Liu~Y.\,W., Fang~Q.\,H.

\paper Interaction of an anti-plane singularity with interfacial anti-cracks in cylindrically anisotropic composites 

\jour Arch. Appl. Mech.

\vol 78

\issue 4

\pages 295--309

\yr 2008



\RBibitem{bogan:bib4}

\by Прусов~И.\,А.

\book Некоторые задачи термоупругости

\publaddr Минск

\publ БГУ

\yr 1972

\totalpages 200

\misctransl

\by Prusov~I.\,A.

\book Some Problems of Thermoelasticity 

\publaddr Minsk

\publ BGU

\yr 1972

\totalpages 200



\RBibitem{bogan:bib5}

\by Ландис~Е.\,М.

\book Уравнения второго порядка эллиптического и~параболического типов

\publaddr М.

\publ Наука

\yr 1971

\totalpages 288

\misctransl

\by Landis~E.\,M. 

\book Second order equations of elliptic and parabolic type

\publ Nauka

\publaddr Moscow

\yr 1971

\totalpages 288 



\Bibitem{bogan:bib6}

\by Piccinini~L.\,C, Spagnolo~S.

\paper On the H\"{o}lder continuity of solutions of second order elliptic equations in two variables

\jour Ann. Sc. Norm. Super. Pisa, Sci. Fis. Mat., III. Ser.

\vol 26

\issue 2

\yr 1972

\pages 391--402









\end{thebibliography}





\noindent

\hskip6.5cm \thedate \par\noindent

\hskip6.5cm\therevdate









%

%\newpage

%

%\chead[\fancyplain{}{\scriptsize \sl B\,o\,g\,a\,n~Yu.\,A.}]

%{\fancyplain{}{\scriptsize \sl  Antiplane strain of a cylindrically anisotropic elastic

%bar}}



\makeengtitle







 \newpage

%

% \tableofengcontents

%

%\end{document}

